% !TeX root = ../../book.tex
\section{命题逻辑}
\secbegin{secPropositionalLogic}

每个数学证明都是在做出某些\textit{假设}和实现某些\textit{目标}的背景下书写的。

\begin{itemize}
\item \textbf{假设 } 是已知为真的命题，或我们为了证明某事而假定为真的命题。这即包括已经证明的定理、假定读者已有的先验知识，也包括使用`假设'或`假定'等词明确做出的假设。
\item \textbf{目标 } 是我们为了完成结果的证明而试图证明的命题，或者可能只是证明中的一个步骤。
\end{itemize}

我们的假设和目标都会随着我们写下的短语改变。这也许最好用例子来说明。在下面的 \Cref{exAssumptionsGoals} 中，我们将详细检查 \Cref{propDivisibilityIsTransitive} 的证明，以便我们可以看到我们所写的文字如何影响证明中每个阶段的假设和目标。 我们将在任何给定阶段使用表格表明我们的假设和目标 --- 列出的假设仅是那些明确做出的假设；先验知识和先前证明的定理将被隐含。


\begin{example}
\label{exAssumptionsGoals}
\Cref{propDivisibilityIsTransitive} 的陈述如下：
\begin{quote}
设 $a,b,c \in \mathbb{Z}$。如果 $c$ 整除 $b$ 且 $b$ 整除 $a$，那么 $c$ 整除 $a$。
\end{quote}
命题的设置立即给了我们最初的假设和目标：
\begin{center}
\begin{tabular}{C{150pt}|C{150pt}}
\textbf{假设} & \textbf{目标} \\ \hline
$a,b,c \in \mathbb{Z}$ & 如果 $c$ 整除 $b$ 且 $b$ 整除 $a$，那么 $c$ 整除 $a$
\end{tabular}
\end{center}
接下来我们将逐行进行证明，看看我们写的文字对假设和目标有什么影响。

由于我们的目标是 `如果\dots{}那么\dots{}' 形式的表达式，一个合理的想法是从假设 `如果' 语句开始，然后使用该假设来证明`那么' 语句。因此，我们在证明中写的第一句话是：
\begin{quote}
假设 $c$ 整除 $b$ 且 $b$ 整除 $a$。
\end{quote}
更新后的假设和目标如下表。
\begin{center}
\begin{tabular}{C{150pt}|C{150pt}}
\textbf{假设} & \textbf{目标} \\ \hline
$a,b,c \in \mathbb{Z}$ & $c$ 整除 $a$ \\
$c$ 整除 $b$ & \\
$b$ 整除 $a$ & 
\end{tabular}
\end{center}

我们证明的下一步是给出 `$c$ 整除 $b$' 和 `$b$ 整除 $a$' 的定义，让我们可以更进一步。

\begin{quote}
{\color{gray} 假设 $c$ 整除 $b$ 且 $b$ 整除 $a$。} 根据 \Cref{defDivisionPreliminary}，对于某整数 $q$ 和 $r$，可得
\[
b=qc \quad \text{且} \quad a=rb
\]
\end{quote}

这里引入了两个新变量 $q,r$，并且允许我们用它们的定义替换假设 `$c$ 整除 $b$' 和 `$b$ 整除 $a$'。

\begin{center}
\begin{tabular}{C{150pt}|C{150pt}}
\textbf{假设} & \textbf{目标} \\ \hline
$a,b,c,q,r \in \mathbb{Z}$ & $c$ 整除 $a$ \\
$b=qc$ & \\
$a=rb$ &
\end{tabular}
\end{center}

至此，我们几乎已经用尽了我们可以做出的所有假设，因此我们将注意力转向目标 --- 也就是说，我们必须证明 $c$ 整除 $a$。在这一点上，通过分解目标来`逆向工作'非常有用：$c$ 整除 $a$ 意味着什么？根据 \Cref{defDivisionPreliminary}，我们需要证明 $a$ 等于某个整数乘以 $c$ --- 这将反映在下表的假设和目标中。

由于我们现在试图用 $c$ 来表示 $a$，因此可以使用 $a$ 与 $b$ 以及 $b$ 与 $c$ 关联等式来关联 $a$ 与 $c$。

\begin{quote}
{\color{gray} 假设 $c$ 整除 $b$ 且 $b$ 整除 $a$。根据 \Cref{defDivisionPreliminary}，对于某整数 $q$ 和 $r$，可得
\[
b=qc \quad \text{且} \quad a=rb
\]
} 使用第一个等式，我们可以用 $qc$ 替换第二个等式中的 $b$，得
\[
a=r(qc)
\]
\end{quote}

我们现在离目标很近了，如下表所示。

\begin{center}
\begin{tabular}{C{150pt}|C{150pt}}
\textbf{假设} & \textbf{目标} \\ \hline
$a,b,c,q,r \in \mathbb{Z}$ & $a = [\text{某个整数}] \cdot c$ \\
$b=qc$ & \\
$a=rb$ & \\
$a=r(qc)$ & 
\end{tabular}
\end{center}

最后一步是观察目标终于实现了：

\begin{quote}
{\color{gray} 假设 $c$ 整除 $b$ 且 $b$ 整除 $a$。根据 \Cref{defDivisionPreliminary}，对于某整数 $q$ 和 $r$，可得
\[
b=qc \quad \text{and} \quad a=rb
\]
使用第一个等式，我们可以用 $qc$ 替换第二个等式中的 $b$，得}
\[
{\color{gray} a=r(qc)}
\]
因为 $r(qc) = (rq)c$, 且 $rq$ 为整数，
\end{quote}

\begin{center}
\begin{tabular}{C{150pt}|C{150pt}}
\textbf{假设} & \textbf{目标} \\ \hline
$a,b,c,q,r \in \mathbb{Z}$ &  \\
$b=qc$ & \\
$a=rb$ & \\
$a=r(qc)$ & \\
$a=(rq)c$ & \\
$rq \in \mathbb{Z}$ & 
\end{tabular}
\end{center}

至此所有要证明的都证明了，此时重申观点非常有帮助，这样读者就可以对此事有结论性的了解。

\begin{quote}
{\color{gray} 假设 $c$ 整除 $b$ 且 $b$ 整除 $a$。根据 \Cref{defDivisionPreliminary}，对于某整数 $q$ 和 $r$，可得
\[
b=qc \quad \text{and} \quad a=rb
\]
使用第一个等式，我们可以用 $qc$ 替换第二个等式中的 $b$，得
\[
a=r(qc)
\]
因为 $r(qc) = (rq)c$, 且 $rq$ 为整数，} 所以根据 \Cref{defDivisionPreliminary} 得 $c$ 整除 $a$。
\end{quote}
\end{example}

\subsection*{Symbolic logic}

再次考虑我们在 \Cref{propDivisibilityIsTransitive} 中证明的命题（对于给定整数 $a,b,c$）：

\begin{center}
如果 $c$ 整除 $b$ 且 $b$ 整除 $a$，则 $c$ 整除 $a$。
\end{center}

`$c$ 整除 $b$'、`$b$ 整除 $a$' 和 `$c$ 整除 $a$' 这三个陈述本身就是命题，尽管它们都出现在一个更复杂的命题中。我们可以用符号替换这些更简单的命题来检查命题的逻辑结构，这些符号即所谓的\textit{命题变量}。用 $P$ 表示 `$c$ 整除 $b$'，用 $Q$ 表示 `$b$ 整除 $a$'，用 $R$ 表示 `$c$ 整除 $a$'，我们得到：

\begin{center}
如果 $P$ 且 $Q$，则 $R$。
\end{center}

用这种方式分解命题可以清楚地看到一个可行的\textit{假设} $P$ 和 $Q$，然后从这些假设中\textit{推导}出 $R$ --- 这正是我们在证明中所做的，在 \Cref{exAssumptionsGoals} 中，我们非常详细地检查了证明。但重要的是，它表明相同的证明策略可能适用于其他形式为 `如果 $P$ 且 $Q$，则 $R$' 的命题，例如以下命题（对于给定整数 $n$ ):

\begin{center}
如果 $n > 2$ 且 $n$ 是质数，则 $n$ 是奇数。
\end{center}

通过观察可以发现，简单的命题可以通过词语 `且' 和结构 `如果\dots{} 那么\dots{}' 等语言连接在一起，形成更复杂的命题 --- 我们将使用\textit{逻辑运算符}象征性地表示这些结构，这将在 \Cref{defLogicalOperator} 中介绍。

进一步放大来看，我们可以使用 \Cref{defDivisionPreliminary} 观察到 `$c$ 整除 $b$' 实际上意味着 `对于某个 $q \in \mathbb{Z}$, $b = qc$'。表达式 `对于某个 $q \in \mathbb{Z}$' 引入了一个新变量 $q$，我们必须在证明中恰当地处理它。 我们在证明中附加到变量上的词 --- 例如：`任意', `存在', `所有', `某个', `仅有' 和 `唯一' --- 将使用\textit{量词}来象征性地表示，我们将在 \Cref{secVariablesQuantifiers} 中学习。

通过将一个复杂的命题分解为使用逻辑运算符和量词连接在一起的更简单的陈述，我们可以更加准确地确定，在命题证明的任何给定阶段，我们可以做出哪些假设，以及需要哪些步骤才能完成证明。

\subsection*{命题公式}

我们从固定术语的定义入手发展符号逻辑。

\begin{definition}
\label{defPropositionalVariable}
\index{propositional variable}
\index{truth value}
\textbf{命题变量}是表示命题的符号。命题变量可以赋予\textbf{真值}（`真' 或 `假'）。
\end{definition}

我们通常使用小写字母 $p$、$q$、$r$ 和 $s$ 表示命题变量。

通过使用逻辑运算符将更简单的表达式连接在一起，可以形成更复杂命题的表达式，例如$\wedge$（表示 `与'）、$\vee$（表示`或'）、$\Rightarrow$（表示`如果\dots{}那么\dots{}'）和$\neg $（代表`否'）。

下面给出的\textit{逻辑运算符}和\textit{命题公式}的概念定义有点难理解，因此最好结合命题公式和逻辑运算符的例子来理解。幸运的是，我们会看到很多逻辑运算符和命题公式，因为它们是本节的主要研究对象。

\begin{definition}
\label{defPropositionalFormula}
\label{defLogicalOperator}
\index{logical operator}
\index{propositional formula}
\index{subformula}
\textbf{命题公式}是一个表达式，它要么是一个命题变量，要么是使用 \textbf{\mbox{逻辑} \mbox{运算符}} 从更简单的命题公式（`子公式'）构建的。在后一种情况下，命题公式的真值由子公式的真值根据逻辑运算规则决定。
\end{definition}

乍一看，\Cref{defPropositionalFormula} 似乎是循环定义 --- 用命题公式来定义术语`命题公式'！但实际上这不是循环定义；这是一个\textit{递归}定义的例子（我们用 `更简单' 这个词来避免循环）。为了说明这一点，请考虑以下命题公式示例：
\[
(p \wedge q) \Rightarrow r
\]
这个表达式表示一个形式为 `如果 $p$ 且 $q$，则 $r$' 的命题，其中 $p,q,r$ 本身就是命题。它是使用逻辑运算符 $\Rightarrow$ 从子公式 $p \wedge q$ 和 $r$ 构建的，$p \wedge q$ 本身是使用逻辑运算符 $\wedge$ 从子公式 $p$ 和 $q$ 构建的。

然后 $(p \wedge q) \Rightarrow r$ 的真值，根据逻辑运算符 $\wedge$ 和 $\Rightarrow$ 的规则，由构成命题的变量（$p$、$q$ 和 $r$）的真值决定。

如果这一切看起来有点抽象，那是因为它\textit{就是}抽象。如果你看不懂，这不是你的问题。从现在开始，我们将只研究逻辑运算符的特定实例，这将更容易理解。

\subsubsection*{合取 (`和/且/与', $\wedge$)}

合取是逻辑运算符，当我们说 `和/且/与' 时，它可以准确表达我们的意思。

\begin{idefinition}
\label{defConjunction}
\index{conjunction}
\nindex{conjunction}{$\wedge$}{conjunction}
\textbf{合取}运算符是逻辑运算符 $\wedge$ \inlatex{wedge}\lindexmmc{wedge}{$\wedge$}，根据以下规则定义：
\begin{itemize}
\item \introrule{\wedge} 如果 $p$ 为真且 $q$ 为真，则 $p \wedge q$ 为真；
\item \elimrulesub{\wedge}{1} 如果 $p \wedge q$ 为真，则 $p$ 为真；
\item \elimrulesub{\wedge}{2} 如果 $p \wedge q$ 为真, 则 $q$ 为真。
\end{itemize}
表达式 $p \wedge q$ 表示 `$p$ 且 $q$'。
\end{idefinition}

使用合取时并不总是很明显；有时它会在没有提到 `和/且/与' 这些词的情况下偷偷溜进来！请注意此类情况，例如下面的练习。

\begin{example}
\label{exSevenDividesTwentyEightConjunction}
我们可以用 $p \wedge q$ 这种形式表达 `$7$ 是 $28$ 的质因数' 这一命题，让 $p$ 表示命题 `$7$ 是质数'，让 $q$ 表示命题`$7 $ 整除 $28$'。
\end{example}

\begin{exercise}
\label{exJohnMathematicianPittsburgh}
用 $p \wedge q$ 的形式表达 `约翰是一位生活在匹兹堡的数学家' 这一命题。
\end{exercise}
\Cref{defConjunction} 中的规则是\textit{推理规则}的例子 --- 它们告诉我们如何从其他命题公式的真值中推导（或`推断'）出一个命题公式的真值。特别是，推理规则从不直接告诉我们命题何时为\textit{假} --- 为了证明某事为假，我们将证明其\textit{逆命题}为真（参见 \Cref{defNegation}）。

推理规则告诉我们如何在证明中使用命题的逻辑结构：

\begin{itemize}
\item 规则 \introrule{\wedge} 是一个引入规则，这意味着它告诉我们如何\textit{证明} $p \wedge q$ --- 事实上，如果我们想要证明 $p \wedge q$ 为真，\introrule{\wedge} 告诉我们只要证明 $p$ 为真并且 $q$ 为真就够了。

\item 规则 \elimrulesub{\wedge}{1} 和 \elimrulesub{\wedge}{2} 是消除规则，这意味着它们告诉我们如何使用 $p \wedge q$ 形式的假设 --- 实际上，如果我们假设 $p \wedge q$ 为真，然后我们可以自由使用 $p$ 为真和 $q$ 为真的事实。
\end{itemize}

每个逻辑运算符都会配备一些引入和/或消除规则，这些规则告诉我们如何证明目标或如何使用包含在逻辑运算符中的假设。正是通过这种方式，命题的逻辑结构为\textit{证明策略}提供了信息，如下所示：

\begin{strategy}[合取证明]
\label{strProvingConjunctionsDirect}
命题 $p \wedge q$ 的证明可以通过将两个证明结合在一起获得，一个证明 $p$ 为真，另一个证明 $q$ 为真。
\end{strategy}

\begin{example}
\label{exSevenDividesTwentyEightConjunctionProof}
假设我们需要证明 $7$ 是 $28$ 的质因数。在 \Cref{exSevenDividesTwentyEightConjunction} 中，我们将 `$7$ 是 $28$ 的质因数' 表达为命题 `$7$ 是质数' 和 `$7$ 除 $28$' 的合取，因此 \Cref{strProvingConjunctionsDirect} 将证明分为两步：先证明 $7$ 是素数，再证明 $7$ 能整除 $28$。
\end{example}

就像 \Cref{strProvingConjunctionsDirect} 被 $\wedge$ 的引入规则告知一样，消除规则告诉我们如何使用涉及合取的假设。

\begin{strategy}[合取假设]
\label{strAssumingConjunctionsDirect}
如果证明中的假设具有 $p \wedge q$ 的形式，那么我们可以在证明中假设 $p$ 和假设 $q$。
\end{strategy}

\begin{example}
\label{exSevenDividesTwentyEightConjunctionAssumption}
假设在证明命题的过程中，我们得出这样的事实：$7$ 是 $28$ 的质因数。则 \Cref{strAssumingConjunctionsDirect} 允许我们使用 $7$ 是质数和 $7$ 整除 $28$ 这两个单独的事实。
\end{example}

\Cref{strProvingConjunctionsDirect,strAssumingConjunctionsDirect} 看上去非常\textit{明显}。某种程度上它们是显而易见的，这就是我们首先介绍它们的原因。但我们精确定义逻辑运算符和它们的引入规则与消除规则以及相应的证明策略的真正原因是，当你在证明一个复杂的结果时，很容易忘记你已经证明的和还有待证明的。跟踪证明中的假设和目标，并理解为了完成证明必须做什么，是一项艰巨的任务。

为了避免把这个过程拉得太长，我们需要一种简洁的方式来表达推理规则，让我们能够轻而易举地读出相应的证明策略。我们可以像 \Cref{exAssumptionsGoals} 那样使用假设和目标表，但这很快就会变得笨拙 --- 事实上，即使是非常简单的合取引入规则 \introrule{\wedge} 在这种格式下看起来也不是很好：

\begin{center}
\begin{tabular}{C{60pt}|C{60pt}}
{\small \textbf{假设}} & {\small \textbf{目标}} \\ \hline
$\vdots$ & $p \wedge q$ \\
$\vdots$ & ~
\end{tabular}
$\quad \leadsto \quad$
\begin{tabular}{C{60pt}|C{60pt}}
{\small \textbf{假设}} & \small{\textbf{目标}} \\ \hline
$\vdots$ & $p$ \\
$\vdots$ & $q$
\end{tabular}
\end{center}

相反，我们将以\textit{自然演绎}的方式表示推理规则。在这种风格中，我们将规则的\textit{前提} $p_1,p_2,\dots,p_k$ 写在一条线的上方，在线的下方有一个\textit{结论} $q$，表示断言命题 $q$ 的真实性来自 （所有）前提 $p_1,p_2,\dots,p_k$ 的真实性。

\begin{center}
\begin{prooftree}
  \AxiomC{$p_1$}
  \AxiomC{$p_2$}
  \AxiomC{$\cdots$}
  \AxiomC{$p_k$}
\QuaternaryInfC{$q$}
\end{prooftree}
\end{center}

例如，合取的引入和排除规则可以简明地表示如下：

\begin{center}
\begin{minipage}{0.15\textwidth}
\centering
\begin{prooftree}
  \AxiomC{$p$}
  \AxiomC{$q$}
\TagC{\introrule{\wedge}}
\BinaryInfC{$p \wedge q$}
\end{prooftree}
\end{minipage}
%
\hspace{20pt}
%
\begin{minipage}{0.15\textwidth}
\centering
\begin{prooftree}
  \AxiomC{$p \wedge q$}
\TagC{\elimrulesub{\wedge}{1}}
\UnaryInfC{$p$}
\end{prooftree}
\end{minipage}
%
\hspace{20pt}
%
\begin{minipage}{0.15\textwidth}
\centering
\begin{prooftree}
  \AxiomC{$p \wedge q$}
\TagC{\elimrulesub{\wedge}{2}}
\UnaryInfC{$q$}
\end{prooftree}
\end{minipage}
\end{center}

除了简洁紧凑外，这种编写推理规则的方式还很有用，因为我们可以将它们组合成\textit{证明树}，以便了解如何证明更复杂的命题。例如，考虑下面的证明树，它结合了合取引入规则的两个实例。

\begin{center}
\begin{prooftree}
    \AxiomC{$p$}
    \AxiomC{$q$}
  \BinaryInfC{$p \wedge q$}
  \AxiomC{$r$}
\BinaryInfC{$(p \wedge q) \wedge r$}
\end{prooftree}
\end{center}

从这个证明树中，我们获得了一个证明形式为 $(p \wedge q) \wedge r$ 的命题的策略。即先证明 $p$ 再证明 $q$，得出$p\wedge q$；然后证明 $r$，得出 $(p \wedge q) \wedge r$。这反映出命题的逻辑结构告诉我们如何构建命题的证明。

\begin{exercise}
画一个证明树，其结论是命题公式 $(p \wedge q) \wedge (r \wedge s)$，其中 $p,q,r,s$ 是命题变量。用适当的命题 $p$, $q$, $ r$ 和 $s$，以 $(p \wedge q) \wedge (r \wedge s)$ 的形式表达 `$2$ 是偶素数，$3$ 是奇素数'，并描述你的证明树如何暗示证明可能的样子。
\end{exercise}

\subsubsection*{析取 (`或', $\vee$)}

\begin{idefinition}
\label{defDisjunction}
\index{disjunction}
\nindex{disjunction}{$\vee$}{disjunction}
The \textbf{disjunction} operator is the logical operator $\vee$ \inlatex{vee}\lindexmmc{vee}{$\vee$}, defined according to the following rules:
\begin{itemize}
\item \introrulesub{\vee}{1} If $p$ is true, then $p \vee q$ is true;
\item \introrulesub{\vee}{2} If $q$ is true, then $p \vee q$ is true;
\item \elimrule{\vee} If $p \vee q$ is true, and if $r$ can be derived from $p$ and from $q$, then $r$ is true.
\end{itemize}
The expression $p \vee q$ represents `$p$ or $q$'.
\end{idefinition}

The introduction and elimination rules for disjunction are represented diagramatically as follows.

\begin{center}
\begin{minipage}[b]{0.15\textwidth}
\centering
\begin{prooftree}
  \AxiomC{$p$}
\TagC{\introrulesub{\vee}{1}}
\UnaryInfC{$p \vee q$}
\end{prooftree}
\end{minipage}
%
\hspace{20pt}
%
\begin{minipage}[b]{0.15\textwidth}
\centering
\begin{prooftree}
  \AxiomC{$q$}
\TagC{\introrulesub{\vee}{2}}
\UnaryInfC{$p \vee q$}
\end{prooftree}
\end{minipage}
%
\hspace{20pt}
%
\begin{minipage}[b]{0.35\textwidth}
\begin{prooftree}
  \AxiomC{$p \vee q$}
    \AxiomC{$[p]$}
    \noLine
    \UnaryInfC{$\downleadsto$}
  \noLine
  \UnaryInfC{$r$}
    \AxiomC{$[q]$}
    \noLine
    \UnaryInfC{$\downleadsto$}
  \noLine
  \UnaryInfC{$r$}
\TagC{\elimrule{\vee}}
\TrinaryInfC{$r$}
\end{prooftree}
\end{minipage}
\end{center}

We will discuss what the notation $[p] \leadsto r$ and $[q] \leadsto r$ means momentarily. First, we zoom in on how the disjunction introduction rules inform proofs of propositions of the form `$p$ or $q$'.

\begin{strategy}[Proving disjunctions]
\label{strProvingDisjunctionsDirect}
In order to prove a proposition of the form $p \vee q$, it suffices to prove just one of $p$ or $q$.
\end{strategy}

\begin{example}
Suppose we want prove that $8192$ is not divisible by $3$. We know by the division theorem (\Cref{thmDivisionPreliminary}) that an integer is not divisible by $3$ if and only if it leaves a remainder of $1$ or $2$ when divided by $3$, and so it suffices to prove the following:
\[
\begin{matrix} 8192 \text{ leaves a remainder of } 1 \\ \text{when divided by } 3 \end{matrix}
\quad \vee \quad
\begin{matrix} 8192 \text{ leaves a remainder of } 2 \\
\text{when divided by } 3 \end{matrix}
\]
A quick computation reveals that $8192 = 2730 \times 3 + 2$, so that $8192$ leaves a remainder of $2$ when divided by $3$. By \Cref{strProvingDisjunctionsDirect}, the proof is now complete, since the full disjunction follows by \introrulesub{\vee}{2}.
\end{example}

\begin{example}
Let $p,q,r,s$ be propositional variables. The propositional formula $(p \vee q) \wedge (r \vee s)$ represents `$p$ or $q$, and $r$ or $s$'. What follows are two examples of truth trees for this propositional formula.

\begin{center}
\begin{minipage}{0.3\textwidth}
\centering
\begin{prooftree}
    \AxiomC{$p$}
  \TagC{\introrulesub{\vee}{1}}
  \UnaryInfC{$p \vee q$}
    \AxiomC{$r$}
  \TagC{\introrulesub{\vee}{1}}
  \UnaryInfC{$r \vee s$}
\TagC{\introrule{\wedge}}
\BinaryInfC{$(p \vee q) \wedge (r \vee s)$}
\end{prooftree}
\end{minipage}
%
\hspace{20pt}
%
\begin{minipage}{0.3\textwidth}
\centering
\begin{prooftree}
    \AxiomC{$q$}
  \TagC{\introrulesub{\vee}{2}}
  \UnaryInfC{$p \vee q$}
    \AxiomC{$s$}
  \TagC{\introrulesub{\vee}{2}}
  \UnaryInfC{$r \vee s$}
\TagC{\introrule{\wedge}}
\BinaryInfC{$(p \vee q) \wedge (r \vee s)$}
\end{prooftree}
\end{minipage}
\end{center}

The proof tree on the left suggests the following proof strategy for $(p \vee q) \wedge (r \vee s)$. First prove $p$, and deduce $p \vee q$; then prove $r$, and deduce $r \vee s$; and finally deduce $(p \vee q) \wedge (r \vee s)$. The proof tree on the right suggests a different strategy, where $p \vee q$ is deduced by proving $q$ instead of $p$, and $r \vee s$ is deduced by proving $s$ instead of $r$.

Selecting which (if any) of these to use in a proof might depend on what we are trying to prove. For example, for a fixed natural number $n$, let $p$ represent `$n$ is even', let $q$ represent `$n$ is odd', let $r$ represent `$n \ge 2$' and let $s$ represent `$n$ is a perfect square'. Proving $(p \vee q) \wedge (r \vee s)$ when $n=2$ would be most easily done using the left-hand proof tree above, since $p$ and $r$ are evidently true when $n=2$. However, the second proof tree would be more appropriate for proving $(p \vee q) \wedge (r \vee s)$ when $n=1$.
\end{example}

\begin{aside}
If you haven't already mixed up $\wedge$ and $\vee$, you probably will soon, so here's a way of remembering which is which:
\begin{center} \vspace{-10pt} \large \textbf{fish n chips} \end{center}
\vspace{-10pt} If you forget whether it's $\wedge$ or $\vee$ that means `and', just write it in place of the `n' in `fish n chips':
\begin{center} \vspace{-10pt} fish $\wedge$ chips \qquad \qquad fish $\vee$ chips \end{center}
\vspace{-10pt} Clearly the first looks more correct, so $\wedge$ means `and'. If you don't eat fish (or chips), then worry not, as this mnemonic can be modified to accommodate a wide variety of dietary restrictions; for instance `mac n cheese' or `quinoa n kale' or, for the meat lovers, `ribs n brisket'.
\end{aside}

Recall the diagrammatic statement of the disjunction elimination rule:

\begin{center}
\begin{prooftree}
  \AxiomC{$p \vee q$}
    \AxiomC{$[p]$}
    \noLine
    \UnaryInfC{$\downleadsto$}
  \noLine
  \UnaryInfC{$r$}
    \AxiomC{$[q]$}
    \noLine
    \UnaryInfC{$\downleadsto$}
  \noLine
  \UnaryInfC{$r$}
\TagC{\elimrule{\vee}}
\TrinaryInfC{$r$}
\end{prooftree}
\end{center}

The curious notation $[p] \leadsto r$ indicates that $p$ is a \textit{temporary assumption}. In the part of the proof corresponding to $[p] \leadsto r$, we would assume that $p$ is true and derive $r$ from that assumption, and remove the assumption that $p$ is true from that point onwards. Likewise for $[q] \leadsto r$.

The proof strategy obtained from the disjunction elimination rule is called \textit{proof by cases}.

\begin{strategy}[Assuming disjunctions---proof by cases]
\label{strAssumingDisjunctionsDirect}
\index{proof!by cases}
If an assumption in a proof has the form $p \vee q$, then we may derive a proposition $r$ by splitting into two cases: first, derive $r$ from the temporary assumption that $p$ is true, and then derive $r$ from the assumption that $q$ is true.
\end{strategy}

The following example illustrates how \Cref{strProvingDisjunctionsDirect,strAssumingDisjunctionsDirect} can be used together in a proof.

\begin{example}
\label{exPositiveProperFactorOfFourEvenOrPerfectSquare}
Let $n$ be a positive proper factor of $4$, and suppose we want to prove that $n$ is either even or a perfect square.
\begin{itemize}
\item Our assumption that $n$ is a positive proper factor of $4$ can be expressed as the disjunction $n = 1 \vee n = 2$.
\item Our goal is to prove the disjunction `$n \text{ is even} \vee n \text{ is a perfect square}$'.
\end{itemize}

According to \Cref{strAssumingDisjunctionsDirect}, we split into two cases, one in which $n=1$ and one in which $n=2$. In each case, we must derive `$n \text{ is even} \vee n \text{ is a perfect square}$', for which it suffices by \Cref{strProvingDisjunctionsDirect} to derive either that $n$ is even or that $n$ is a perfect square. Thus a proof might look something like this:

\begin{quote}
Since $n$ is a positive proper factor of $4$, either $n=1$ or $n=2$.
\begin{itemize}
\item \textbf{Case 1.} Suppose $n=1$. Then since $1^2=1$ we have $n = 1^2$, so that $n$ is a perfect square.
\item \textbf{Case 2.} Suppose $n=2$. Then since $2 = 2 \times 1$, we have that $n$ is even.
\end{itemize}
Hence $n$ is either even or a perfect square. \qed
\end{quote}

Notice that in both Case 1 and Case 2, we did not explicitly mention that we had proved that `$n \text{ is even} \vee n \text{ is a perfect square}$', leaving that deduction to the reader---we only mentioned it after the proofs in each case were complete.
\end{example}

The proof of \Cref{propRemainderOfSquaresModulo3} below splits into \textit{three} cases, rather than just two.

\begin{proposition}
\label{propRemainderOfSquaresModulo3}
Let $n \in \mathbb{Z}$. Then $n^2$ leaves a remainder of $0$ or $1$ when divided by $3$.
\end{proposition}

\begin{cproof}
Let $n \in \mathbb{Z}$. By the division theorem (\Cref{thmDivisionPreliminary}), one of the following must be true for some $k \in \mathbb{Z}$:
\[
n=3k \quad \text{or} \quad n=3k+1 \quad \text{or} \quad n=3k+2
\]
\begin{itemize}
\item Suppose $n=3k$. Then
\[
n^2 = (3k)^2 = 9k^2 = 3 \cdot (3k^2)
\]
So $n^2$ leaves a remainder of $0$ when divided by $3$.
\item Suppose $n=3k+1$. Then
\[
n^2 = (3k+1)^2 = 9k^2+6k+1 = 3(3k^2+2k)+1
\]
So $n^2$ leaves a remainder of $1$ when divided by $3$.
\item Suppose $n=3k+2$. Then
\[
n^2 = (3k+2)^2 = 9k^2+12k+4 = 3(3k^2+4k+1)+1
\]
So $n^2$ leaves a remainder of $1$ when divided by $3$.
\end{itemize}
In all cases, $n^2$ leaves a remainder of $0$ or $1$ when divided by $3$.
\end{cproof}

Note that in the proof of \Cref{propRemainderOfSquaresModulo3}, unlike in \Cref{exPositiveProperFactorOfFourEvenOrPerfectSquare}, we did not explictly use the word `case', even though we were using proof by cases. Whether or not to make your proof strategies explicit is up to you---discussion of this kind of matter can be found in \Cref{secVocabulary}.

When completing the following exercises, try to keep track of exactly where you use the introduction and elimination rules that we have seen so far.

\begin{exercise}
Let $n$ be an integer. Prove that $n^2$ leaves a remainder of $0$, $1$ or $4$ when divided by $5$.
\hintlabel{exSquareRemainderModuloFive}{f
Mimic the proof of \Cref{propRemainderOfSquaresModulo3}.
}
\end{exercise}

\begin{exercise}
Let $a,b \in \mathbb{R}$ and suppose $a^2-4b \ne 0$. Let $\alpha$ and $\beta$ be the (distinct) roots of the polyonomial $x^2+ax+b$. Prove that there is a real number $c$ such that either $\alpha-\beta = c$ or $\alpha - \beta = ci$.
\end{exercise}

\subsubsection*{Implication (`if\dots{}then\dots{}', $\Rightarrow$)}

\begin{definition}
\label{defImplication}
\index{implication}
\nindex{implies}{$\Rightarrow$}{implication}
The \textbf{implication} operator is the logical operator $\Rightarrow$ \inlatex{Rightarrow}\lindexmmc{Rightarrow}{$\Rightarrow$}, defined according to the following rules:
\begin{itemize}
\item \introrule{\Rightarrow} If $q$ can be derived from the assumption that $p$ is true, then $p \Rightarrow q$ is true;
\item \elimrule{\Rightarrow} If $p \Rightarrow q$ is true and $p$ is true, then $q$ is true.
\end{itemize}
The expression $p \Rightarrow q$ represents `if $p$, then $q$'.
\end{definition}

\begin{center}
\begin{minipage}[b]{0.15\textwidth}
\begin{prooftree}
      \AxiomC{$[p]$}
    \noLine
    \UnaryInfC{$\downleadsto$}
  \noLine
  \UnaryInfC{$q$}
\TagC{\introrule{\Rightarrow}}
\UnaryInfC{$p \Rightarrow q$}
\end{prooftree}
\end{minipage}
%
\vspace{20pt}
%
\begin{minipage}[b]{0.3\textwidth}
\begin{prooftree}
  \AxiomC{$p \Rightarrow q$}
  \AxiomC{$p$}
\TagC{\elimrule{\Rightarrow}}
\BinaryInfC{$q$}
\end{prooftree}
\end{minipage}
\end{center}

\begin{strategy}[Proving implications]
\label{strProvingImplicationsDirect}
In order to prove a proposition of the form $p \Rightarrow q$, it suffices to assume that $p$ is true, and then derive $q$ from that assumption.
\end{strategy}

The following proposition illustrates how \Cref{strProvingImplicationsDirect} can be used in a proof.

\begin{proposition}
\label{propRationalTwoOfThree}
Let $x$ and $y$ be real numbers. If $x$ and $x+y$ are rational, then $y$ is rational.
\end{proposition}

\begin{cproof}
Suppose $x$ and $x+y$ are rational. Then there exist integers $a,b,c,d$ with $b,d \ne 0$ such that
\[
x = \frac{a}{b} \quad \text{and} \quad x+y = \frac{c}{d}
\]
It then follows that
\[
y = (x+y)-x = \frac{c}{d}-\frac{a}{b} = \frac{bc-ad}{bd}
\]
Since $bc-ad$ and $bd$ are integers, and $bd \ne 0$, it follows that $y$ is rational.
\end{cproof}

The key phrase in the above proof was `Suppose $x$ and $x+y$ are rational.' This introduced the assumptions $x \in \mathbb{Q}$ and $x+y \in \mathbb{Q}$, and reduced our goal to that of deriving a proof that $y$ is rational---this was the content of the rest of the proof.

\begin{exercise}
Let $p(x)$ be a polynomial over $\mathbb{C}$. Prove that if $\alpha$ is a root of $p(x)$, and $a \in \mathbb{C}$, then $\alpha$ is a root of $(x-a)p(x)$.
\end{exercise}

The elimination rule for implication \elimrule{\Rightarrow} is more commonly known by the Latin name \textit{modus ponens}.

\begin{strategy}[Assuming implications---modus ponens]
\label{strAssumingImplicationsDirect}
\index{modus ponens}
If an assumption in a proof has the form $p \Rightarrow q$, and $p$ is also assumed to be true, then we may deduce that $q$ is true.
\end{strategy}

\Cref{strAssumingDisjunctionsDirect} is frequently used to reduce a more complicated goal to a simpler one. Indeed, if we know that $p \Rightarrow q$ is true, and if $p$ is easy to verify, then it allows us to prove $q$ by proving $p$ instead.

\begin{example}
Let $f(x) = x^2+ax+b$ be a polynomial with $a,b \in \mathbb{R}$, and let $\Delta = a^2-4b$ be its discriminant. Part of \Cref{exDiscriminantRealRoots} was to prove that:
\begin{enumerate}[(i)]
\item If $\Delta > 0$, then $f$ has two real roots;
\item If $\Delta = 0$, then $f$ has one real root;
\item If $\Delta < 0$, then $f$ has no real roots.
\end{enumerate}
Given the polynomial $f(x) = x^2-68x+1156$, it would be a pain to go through the process of solving the equation $f(x)=0$ in order to determine how many real roots $f$ has. However, each of the propositions (i), (ii) and (iii) take the form $p \Rightarrow q$, so \Cref{strAssumingImplicationsDirect} reduces the problem of finding how many real roots $f$ has to that of evaluating $\Delta$ and comparing it with $0$. And indeed, $(-68)^2 - 4 \times 1156 = 0$, so the implication (ii) together with \elimrule{\Rightarrow} tell us that $f$ has one real root.
\end{example}

A common task faced by mathematicians is to prove that two conditions are equivalent. For example, given a polynomial $f(x) = x^2+ax+b$ with $a,b \in \mathbb{R}$, we know that \textit{if} $a^2-4b>0$ \textit{then} $f$ has two real roots, but is it also true that if $f$ has two real roots then $a^2-4b>0$? (The answer is `yes'.) The relationship between these two implications is that each is the \textit{converse} of the other.

\begin{definition}
\label{defConverse}
\index{converse}
The \textbf{converse} of a proposition of the form $p \Rightarrow q$ is the proposition $q \Rightarrow p$.
\end{definition}

A quick remark on terminology is pertinent. The following table summarises some common ways of referring to the propositions `$p \Rightarrow q$' and `$q \Rightarrow p$'.

\begin{center}
\begin{tabular}{c|c}
$p \Rightarrow q$ & $q \Rightarrow p$ \\ \hline
if $p$, then $q$ & if $q$, then $p$ \\
$p$ only if $q$ & $p$ if $q$ \\
$p$ is sufficient for $q$ & $p$ is necessary for $q$
\end{tabular}
\end{center}

We so often encounter the problem of proving both an implication and its converse that we introduce a new logical operator that represents the conjunction of both.

\begin{definition}
\label{defBiconditional}
\index{biconditional}
\nindex{biconditional}{$\Leftrightarrow$}{biconditional}
The \textbf{biconditional} operator is the logical operator $\Leftrightarrow$ \inlatex{Leftrightarrow}\lindexmmc{Leftrightarrow}{$\Leftrightarrow$}, defined by declaring $p \Leftrightarrow q$ to mean $(p \Rightarrow q) \wedge (q \Rightarrow p)$. The expression $p \Leftrightarrow q$ represents `$p$ if and only if $q$'.
\end{definition}

Many examples of biconditional statements come from solving equations; indeed, to say that the values $\alpha_1,\dots,\alpha_n$ are the solutions to a particular equation is precisely to say that
\[
x \text{ is a solution} \quad \Leftrightarrow \quad x = \alpha_1 \text{ or } x = \alpha_2 \text{ or } \cdots \text{ or } x = \alpha_n
\]

\begin{example}
\label{exSolveSqrtFirstExample}
We find all real solutions $x$ to the equation
\[
\sqrt{x-3} + \sqrt{x+4} = 7
\]
Let's rearrange the equation to find out what the possible solutions may be.
\begin{align*}
&\phantom{\Rightarrow\;\;} \sqrt{x-3} + \sqrt{x+4} = 7 && \\
&\Rightarrow (x-3) + 2\sqrt{(x-3)(x+4)} + (x+4) = 49 && \text{squaring} \\
&\Rightarrow 2\sqrt{(x-3)(x+4)} = 48-2x && \text{rearranging} \\
&\Rightarrow 4(x-3)(x+4) = (48-2x)^2 && \text{squaring} \\
&\Rightarrow 4x^2+4x-48 = 2304-192x+4x^2 && \text{expanding} \\
&\Rightarrow 196x = 2352 && \text{rearranging} \\
&\Rightarrow x=12 && \text{dividing by $196$}
\end{align*}
You might be inclined to stop here. Unfortunately, all we have proved is that, given a real number $x$, \textit{if} $x$ solves the equation $\sqrt{x-3} + \sqrt{x+4} = 7$, \textit{then} $x=12$. This narrows down the set of possible solutions to just one candidate---but we still need to check the converse, namely that \textit{if} $x=12$, \textit{then} $x$ is a solution to the equation.

As such, to finish off the proof, note that
\[
\sqrt{12-3} + \sqrt{12+4} = \sqrt{9} + \sqrt{16} = 3 + 4 = 7
\]
and so the value $x=12$ is indeed a solution to the equation.
\end{example}

The last step in \Cref{exSolveSqrtFirstExample} may have seemed a little bit silly; but \Cref{exSolveSqrtSecondExample} demonstrates that proving the converse when solving equations truly is necessary.

\begin{example}
\label{exSolveSqrtSecondExample}
We find all real solutions $x$ to the equation
\[
x+\sqrt{x}=0
\]
We proceed as before, rearranging the equation to find all possible solutions.
\begin{align*}
&\phantom{\Rightarrow\;\;} x+\sqrt{x} = 0 && \\
&\Rightarrow x=-\sqrt{x} && \text{rearranging} \\
&\Rightarrow x^2=x && \text{squaring} \\
&\Rightarrow x(x-1)=0 && \text{rearranging} \\
&\Rightarrow x=0 \text{ or } x=1 && 
\end{align*}
Now certainly $0$ is a solution to the equation, since
\[
0+\sqrt{0} = 0+0 = 0
\]
However, $1$ is \textit{not} a solution, since
\[
1+\sqrt{1} = 1+1 = 2
\]
Hence it is actually the case that, given a real number $x$, we have
\[
x+\sqrt{x} = 0 \quad \Leftrightarrow \quad x=0
\]
Checking the converse here was vital to our success in solving the equation!
\end{example}

A slightly more involved example of a biconditional statement arising from the solution to an equation---in fact, a class of equations---is the proof of the quadratic formula.

\begin{itheorem}[Quadratic formula]
\label{thmQuadraticFormula}
Let $a,b \in \mathbb{C}$. A complex number $\alpha$ is a root of the polynomial $x^2+ax+b$ if and only if
\[
\alpha = \frac{-a+\sqrt{a^2-4b}}{2} \quad \text{or} \quad \alpha =\frac{-a-\sqrt{a^2-4b}}{2}
\]
\end{itheorem}

\begin{cproof}
First we prove that \textit{if} $\alpha$ is a root, \textit{then} $\alpha$ is one of the values given in the statement of the proposition. So suppose $\alpha$ be a root of the polynomial $x^2+ax+b$. Then
\[
\alpha^2 + a\alpha + b = 0
\]
The algebraic technique of `completing the square' tells us that
\[
\alpha^2 + a\alpha = \left( \alpha + \frac{a}{2} \right)^2 - \frac{a^2}{4}
\]
and hence
\[
\left( \alpha + \frac{a}{2} \right)^2 - \frac{a^2}{4} + b = 0
\]
Rearranging yields
\[
\left( \alpha + \frac{a}{2} \right)^2  = \frac{a^2}{4} - b = \frac{a^2-4b}{4}
\]
Taking square roots gives
\[
\alpha + \frac{a}{2} = \frac{\sqrt{a^2-4b}}{2} \quad \text{or} \quad \alpha + \frac{a}{2} = \frac{-\sqrt{a^2-4b}}{2}
\]
and, finally, subtracting $\frac{a}{2}$ from both sides gives the desired result.

The proof of the converse is \Cref{exQuadraticFormulaConverse}.
\end{cproof}

\begin{exercise}
\label{exQuadraticFormulaConverse}
Complete the proof of the quadratic formula. That is, for fixed $a,b \in \mathbb{C}$, prove that if
\[
\alpha = \frac{-a+\sqrt{a^2-4b}}{2} \quad \text{or} \quad \alpha =\frac{-a-\sqrt{a^2-4b}}{2}
\]
then $\alpha$ is a root of the polynomial $x^2+ax+b$.
\end{exercise}

Another class of examples of biconditional propositions arise in finding necessary and sufficient criteria for an integer $n$ to be divisible by some number---for example, that an integer is divisible by $10$ if and only if its base-$10$ expansion ends with the digit $0$.

\begin{example}
\label{exTestForDivisibilityByEight}
Let $n \in \mathbb{N}$. We will prove that $n$ is divisible by $8$ if and only if the number formed of the last three digits of the base-$10$ expansion of $n$ is divisible by $8$.

First, we will do some `scratch work'. Let $d_rd_{r-1}\dots{}d_1d_0$ be the base-$10$ expansion of $n$. Then
\[
n = d_r \cdot 10^r + d_{r-1} \cdot 10^{r-1} + \cdots + d_1 \cdot 10 + d_0
\]
Define
\[
n' = d_2d_1d_0 \quad \text{and} \quad n'' = n-n' = d_rd_{r-1}\dots{}d_4d_3000
\]
Now $n-n' = 1000 \cdot d_rd_{r-1} \dots d_4d_3$ and $1000 = 8 \cdot 125$, so it follows that $8$ divides $n''$.

%% BEGIN EXTRACT (xtrStepsExample) %%
Our goal is now to prove that $8$ divides $n$ if and only if $8$ divides $n'$.

\begin{itemize}
\item ($\Rightarrow$) Suppose $8$ divides $n$. Since $8$ divides $n''$, it follows from \Cref{exDivisibilityIsLinear} that $8$ divides $an+bn''$ for all $a,b \in \mathbb{Z}$. But
\[
n' = n-(n-n') = n-n'' = 1 \cdot n + (-1) \cdot n''
\]
so indeed $8$ divides $n'$, as required.
\item ($\Leftarrow$)  Suppose $8$ divides $n'$. Since $8$ divides $n''$, it follows from \Cref{exDivisibilityIsLinear} that $8$ divides $an'+bn''$ for all $a,b \in \mathbb{Z}$. But
\[
n = n'+(n-n') = n'+n'' = 1 \cdot n' + 1 \cdot n''
\]
so indeed $8$ divides $n$, as required.
\end{itemize}
%% END EXTRACT
\end{example}

\begin{exercise}
Prove that a natural number $n$ is divisible by $3$ if and only if the sum of its base-$10$ digits is divisible by $3$.
\hintlabel{exSumOfDigitsDivisibleByThree}{%
Suppose $n = d_r \cdot 10^r + \cdots + d_1 \cdot 10 + d_0$ and let $s = d_r + \cdots + d_1 + d_0$. Start by proving that $3 \mid n-s$.
}
\end{exercise}

\subsubsection*{Negation (`not', $\neg$)}

So far we only officially know how to prove that true propositions are \textit{true}. The negation operator makes precise what we mean by `not', which allows us to prove that false propositions are \textit{false}.

\begin{definition}
\label{defContradiction}
\index{contradiction}
\nindex{contradiction}{$\bot$}{contradiction}
A \textbf{contradiction} is a proposition that is known or assumed to be false. We will use the symbol $\bot$ \inlatex{bot}\lindexmmc{bot}{$\bot$} to represent an arbitrary contradiction.
\end{definition}

\begin{example}
Some examples of contradictions include the assertion that $0=1$, or that $\sqrt{2}$ is rational, or that the equation $x^2=-1$ has a solution $x \in \mathbb{R}$.
\end{example}

\begin{definition}
\label{defNegation}
\index{negation}
\nindex{negation}{$\neg$}{negation}
The \textbf{negation} operator is the logical operator $\neg$ \inlatex{neg}\lindexmmc{neg}{$\neg$}, defined according to the following rules:
\begin{itemize}
\item \introrule{\neg} If a contradiction can be derived from the assumption that $p$ is true, then $\neg p$ is true;
\item \elimrule{\neg} If $\neg p$ and $p$ are both true, then a contradiction may be derived.
\end{itemize}
The expression $\neg p$ represents `not $p$' (or `$p$ is false').
\end{definition}

\begin{center}
\begin{minipage}[b]{0.2\textwidth}
\begin{prooftree}
      \AxiomC{$[p]$}
    \noLine
    \UnaryInfC{$\downleadsto$}
  \noLine
  \UnaryInfC{$\bot$}
\TagC{\introrule{\neg}}
\UnaryInfC{$\neg p$}
\end{prooftree}
\end{minipage}
%
\hspace{20pt}
%
\begin{minipage}[b]{0.2\textwidth}
\begin{prooftree}
  \AxiomC{$\neg p$}
  \AxiomC{$p$}
\TagC{\elimrule{\neg}}
\BinaryInfC{$\bot$}
\end{prooftree}
\end{minipage}
\end{center}

\begin{aside}
The rules \introrule{\neg} and \elimrule{\neg} closely resemble \introrule{\Rightarrow} and \elimrule{\Rightarrow}---indeed, we could simply define $\neg p$ to mean `$p \Rightarrow \bot$', where $\bot$ represents an arbitrary contradiction, but it will be easier later on to have a primitive notion of negation.
\end{aside}

The introduction rule for negation \introrule{\neg} gives rise to a proof strategy called \textit{proof by contradiction}, which turns out to be extremely useful.

\begin{strategy}[Proving negations---proof by contradiction]
\label{strProvingNegationsDirect}
\label{strProofByContradictionDirect}
\index{proof!by contradiction (direct)}
\index{contradiction!(direct) proof by}
In order to prove a proposition $p$ is false (that is, that $\neg p$ is true), it suffices to assume that $p$ is true and derive a contradiction.
\end{strategy}

The following proposition has a classic proof by contradiction.

\begin{proposition}
Let $r$ be a rational number and let $a$ be an irrational number. Then $r+a$ is irrational. 
\end{proposition}

\begin{cproof}
By \Cref{defIrrationalNumber}, we need to prove that $r+a$ is real and not rational. It is certainly real, since $r$ and $a$ are real, so it remains to prove that $r+a$ is not rational.

Suppose $r+a$ is rational. Since $r$ is rational, it follows from \Cref{propRationalTwoOfThree} that $a$ is rational, since
\[
a = (r+a) - r
\]
This contradicts the assumption that $a$ is irrational. It follows that $r+a$ is not rational, and is therefore irrational.
\end{cproof}

Now you can try proving some elementary facts by contradiction.

\begin{exercise}
\label{exNegationAndReciprocalOfIrrationalNumbers}
Let $x \in \mathbb{R}$. Prove by contradiction that if $x$ is irrational then $-x$ and $\frac{1}{x}$ are irrational.
\end{exercise}

\begin{exercise}
\label{exNoLeastPositiveReal}
Prove by contradiction that there is no least positive real number. That is, prove that there is not a positive real number $a$ such that $a \le b$ for all positive real numbers $b$.
\end{exercise}

A proof need not be a `proof by contradiction' in its entirety---indeed, it may be that only a small portion of the proof uses contradiction. This is exhibited in the proof of the following proposition.

\begin{proposition}
\label{propOddIffRemainderOfOne}
Let $a$ be an integer. Then $a$ is odd if and only if $a=2b+1$ for some integer $b$.
\end{proposition}
\begin{cproof}
Suppose $a$ is odd. By the division theorem (\Cref{thmDivisionPreliminary}), either $a=2b$ or $a=2b+1$, for some $b \in \mathbb{Z}$. If $a=2b$, then $2$ divides $a$, contradicting the assumption that $a$ is odd; so it must be the case that $a=2b+1$.

Conversely, suppose $a=2b+1$. Then $a$ leaves a remainder of $1$ when divided by $2$. However, by the division theorem, the even numbers are precisely those that leave a remainder of $0$ when divided by $2$. It follows that $a$ is not even, so is odd.
\end{cproof}

The elimination rule for the negation operator \elimrule{\neg} simply says that a proposition can't be true and false at the same time.

\begin{strategy}[Assuming negations]
\label{strAssumingNegations}
If an assumption in a proof has the form $\neg p$, then any derivation of $p$ leads to a contradiction.
\end{strategy}

The main use of \Cref{strAssumingNegations} is for obtaining the contradiction in a proof by contradiction---in fact, we have already used it in our examples of proof by contradiction! As such, we will not dwell on it further.

\subsection*{Logical axioms}

We wrap up this section by introducing a couple of additional logical rules (\textit{axioms}) that we will use in our proofs.

The first is the so-called \textit{law of excluded middle}, which appears so obvious that it is not even worth stating (let alone naming)---what it says is that every proposition is either true or false. But beware, as looks can be deceiving; the law of excluded middle is a non-constructive axiom, meaning that it should not be accepted in settings it is important to keep track of how a proposition is proved---simply knowing that a proposition is either true or false tells us nothing about how it might be proved or refuted. In most mathematical contexts, though, it is accepted without a second's thought.

\begin{axiom}[Law of excluded middle]
\label{axLEM}
\index{law of excluded middle}
Let $p$ be a propositional formula. Then $p \vee (\neg p)$ is true.
\end{axiom}

The law of excluded middle can be represented diagramatically as follows; there are no premises above the line, since we are simply asserting that it is true.

\begin{center}
\begin{prooftree}
  \AxiomC{}
\TagC{LEM}
\UnaryInfC{$p \vee (\neg p)$}
\end{prooftree}
\end{center}

\begin{strategy}[Using the law of excluded middle]
\label{strLEM}
In order to prove a proposition $q$ is true, it suffices to split into cases based on whether some other proposition $p$ is true or false, and prove that $q$ is true in each case.
\end{strategy}

The proof of \Cref{propIfProductEvenThenSomeFactorEven} below makes use of the law of excluded middle---note that we defined `odd' to mean `not even' (\Cref{defEvenOdd}).

\begin{proposition}
\label{propIfProductEvenThenSomeFactorEven}
Let $a,b \in \mathbb{Z}$. If $ab$ is even, then either $a$ is even or $b$ is even (or both).
\end{proposition}
\begin{cproof}
Suppose $a,b \in \mathbb{Z}$ with $ab$ even.
\begin{itemize}
\item Suppose $a$ is even---then we're done.
\item Suppose $a$ is odd. If $b$ is also odd, then by \Cref{propOddIffRemainderOfOne} can write
\[
a = 2k+1 \quad \text{and} \quad b=2\ell+1
\]
for some integers $k,\ell$. This implies that
\[
ab = (2k+1)(2\ell+1) = 4k\ell + 2k + 2\ell + 1 = 2(\underbrace{2k\ell + k + \ell}_{\in \mathbb{Z}}) + 1
\]
so that $ab$ is odd. This contradicts the assumption that $ab$ is even, and so $b$ must in fact be even.
\end{itemize}
In both cases, either $a$ or $b$ is even.
\end{cproof}

\begin{exercise}
Reflect on the proof of \Cref{propIfProductEvenThenSomeFactorEven}. Where in the proof did we use the law of excluded middle? Where in the proof did we use proof by contradiction? What was the contradiction in this case? Prove \Cref{propIfProductEvenThenSomeFactorEven} twice more, once using contradiction and not using the law of excluded middle, and once using the law of excluded middle and not using contradiction.
\end{exercise}

\begin{exercise}
Let $a$ and $b$ be irrational numbers. By considering the number $\sqrt{2}^{\sqrt{2}}$, prove that it is possible that $a^b$ be rational.
\hintlabel{exIrrationalExpIrrationalCanBeRational}{%
Use the law of excluded middle according to whether the proposition `$\sqrt{2}^{\sqrt{2}}$ is rational' is true or false.}
\end{exercise}

Another logical rule that we will use is the \textit{principle of explosion}, which is also known by its Latin name, \textit{ex falso sequitur quodlibet}, which approximately translates to `\textit{from falsity follows whatever you like}'.

\begin{axiom}[Principle of explosion]
\label{axPrincipleOfExplosion}
If a contradiction is assumed, any consequence may be derived.
\end{axiom}

\begin{center}
\begin{prooftree}
  \AxiomC{$\bot$}
\TagC{Expl}
\UnaryInfC{$p$}
\end{prooftree}
\end{center}

The principle of explosion is a bit confusing on first sight. To shed a tiny bit of intuition on it, think of it as saying that both true and false propositions are consequences of a contradictory assumption. For instance, suppose that $-1 = 1$. From this we can obtain consequences that are false, such as $0=2$ by adding $1$ to both sides of the equation, and consequences that are true, such as $1=1$ by squaring both sides of the equation.

We will rarely use the principle of explosion directly in our mathematical proofs, but we will use it in \Cref{secLogicalEquivalence} for proving logical formulae are equivalent.

\begin{tldr}{secPropositionalLogic}

\subsubsection*{Propositional formulae}

\begin{tldrlist}
\tldritem{defPropositionalVariable}
Propositional variables $p,q,r,\dots$ are used to express simple propositions.

\tldritem{defPropositionalFormula}
Propositional formulae are more complicated expressions built from propositional variables using logical operators, which represent phrases like `and', `or' and `if\dots{} then\dots{}'.
\end{tldrlist}

\subsubsection*{Logical operators}

\begin{tldrlist}
\tldritem{defConjunction}
The \textit{conjunction} operator ($\wedge$) represents `and'. We prove $p \wedge q$ by proving both $p$ and $q$ separately; we can use an assumption of the form $p \wedge q$ by assuming both $p$ and $q$ separately.

\tldritem{defDisjunction}
The \textit{disjunction} operator ($\vee$) represents `or'. We prove $p \vee q$ by proving at least one of $p$ or $q$; we can use an assumption of the form $p \vee q$ by splitting into cases, assuming that $p$ is true in one case, and assuming that $q$ is true in the other.

\tldritem{defImplication}
The \textit{implication} operator ($\Rightarrow$) represents `if\dots{} then\dots{}'. We prove $p \Rightarrow q$ by assuming $p$ and deriving $q$; we can use an assumption of the form $p \Rightarrow q$ by deducing $q$ whenever we know that $p$ is true.

\tldritem{defNegation}
The \textit{negation} operator] ($\neg$) represents `not'. We prove $\neg p$ by assuming $p$ and deriving something known or assumed to be false (this is called \textit{proof by contradiction}); we can use an assumption of the form $\neg p$ by arriving at a contradiction whenever we find that $p$ is true.

\tldritem{defBiconditional}
The \textit{biconditional} operator ($\Leftrightarrow$) represents `if and only if'. The expression $p \Leftrightarrow q$ is shorthand for $(p \Rightarrow q) \wedge (q \Rightarrow p)$.
\end{tldrlist}

\subsubsection*{Logical axioms}

\begin{tldrlist}
\tldritem{axLEM}
The \textit{law of excluded middle} says that $p \vee (\neg p)$ is always true. So at any point in a proof, we may split into two cases: one where $p$ is assumed to be true, and one where $p$ is assumed to be false.

\tldritem{axPrincipleOfExplosion}
The\textit{ principle of explosion} says that if a contradiction is introduced as an assumption, then any conclusion at all can be derived.
\end{tldrlist}
\end{tldr}
